\documentclass[11pt]{article}
\usepackage[margin=0.9in]{geometry}
\usepackage{amsmath,amssymb,amsthm,mathtools,booktabs,longtable,array,microtype}
\usepackage[dvipsnames]{xcolor}
\usepackage[colorlinks=true,linkcolor=NavyBlue,urlcolor=BrickRed]{hyperref}
\usepackage{fancyhdr,enumitem}

\definecolor{ink}{HTML}{17212B}
\definecolor{navy}{HTML}{17365D}
\definecolor{teal}{HTML}{147D78}
\definecolor{pale}{HTML}{EEF4F7}
\color{ink}
\pagestyle{fancy}\fancyhf{}
\lhead{Power characters and bounded support}\rhead{Proof and exact verification}\cfoot{\thepage}
\setlength{\headheight}{14pt}
\setlength{\parindent}{0pt}\setlength{\parskip}{5pt}
\setlist[itemize]{leftmargin=1.35em,itemsep=2pt,topsep=3pt}

\newtheorem{theorem}{Theorem}
\newtheorem{proposition}[theorem]{Proposition}
\newtheorem{remark}[theorem]{Remark}
\newcommand{\Sym}{\operatorname{Sym}}
\newcommand{\Tr}{\operatorname{Tr}}
\newcommand{\Fix}{\operatorname{Fix}}
\newcommand{\ExpS}{\operatorname{Exp}_{\sigma}}
\newcommand{\MGF}{\mathcal S}
\newcommand{\F}{\mathbb F}
\newcommand{\C}{\mathbb C}
\newcommand{\E}{\mathbb E}
\newcommand{\resultbox}[1]{\begin{center}\fcolorbox{navy}{pale}{\parbox{0.93\linewidth}{#1}}\end{center}}

\title{The Power-Character and Bounded-Support Stability Theorem\\
\large Proof, scope, and exact matrix verification}
\author{Lambda Distributions workspace}
\date{20 July 2026}

\begin{document}
\maketitle

\begin{abstract}
For an arbitrary finite group and honest complex representation we derive the
sigma-MGF from character values on all powers.  For permutation actions this
becomes an exact fixed-point and cycle formula; for coset actions it becomes a
centralizer--subgroup-intersection formula, and for conjugation it becomes a
centralizer formula.  We then prove coefficientwise fixed-field stability for
bounded flags in the classical growing-rank towers.  A marimo notebook checks
the linear, parabolic, conjugation, and stability statements by explicit
matrices or generator orbits.  The proof is general; finite computations are
reported only as independent exact verification.
\end{abstract}

\resultbox{\textbf{Status.} Equations \eqref{eq:molien}--\eqref{eq:cosetcycles}
and the type-$A$ threshold \eqref{eq:typeAstable} are exact theorems.  The
formed-space statement is coefficientwise stability with a finite
tower-dependent threshold, not a claimed universal numerical threshold.
Complete flags are excluded: their degree-two coefficient already grows with
the Weyl group.}

\section{Common target and universal character-power formula}

Let $G$ be finite and let $\rho:G\to GL(V)$ be an honest finite-dimensional
complex representation with character $\chi_V$.  For a partition
$\tau=(\tau_1,\ldots,\tau_s)$ set
\[
 W_\tau(V)=\bigotimes_{j=1}^s\Sym^{\tau_j}V,
 \qquad |\tau|=\sum_j\tau_j,
\]
and define
\[
 \MGF_{G,V}=\E_{g\in G}\!\left[\ExpS(X_V(g)h_1)\right]
 =\sum_\tau\dim W_\tau(V)^G\,m_\tau.
\]

\begin{theorem}[Universal power-character formula]\label{thm:power}
One has the exact identities
\begin{align}
\MGF_{G,V}
 &=\frac1{|G|}\sum_{g\in G}\prod_{i\ge1}
       \det(I-t_i\rho(g))^{-1},\label{eq:molien}\\
 &=\frac1{|G|}\sum_{g\in G}\exp\!\left(
       \sum_{r\ge1}\frac{\chi_V(g^r)}r p_r\right),\label{eq:power}\\
 &=\sum_{[g]\subseteq G}\frac1{|C_G(g)|}\exp\!\left(
       \sum_{r\ge1}\frac{\chi_V(g^r)}r p_r\right).\label{eq:classpower}
\end{align}
Thus the series is determined by conjugacy classes, class power maps, and the
character values.  Its coefficient is
\begin{equation}\label{eq:coefficient}
[m_\tau]\MGF_{G,V}=\frac1{|G|}\sum_{g\in G}\prod_{j=1}^s
\left(\sum_{\lambda\vdash\tau_j}\frac1{z_\lambda}
\prod_{r\ge1}\chi_V(g^r)^{m_r(\lambda)}\right).
\end{equation}
\end{theorem}

\begin{proof}
On $W_\tau(V)$ the Reynolds operator
$|G|^{-1}\sum_g\rho_{W_\tau}(g)$ is an idempotent with image
$W_\tau(V)^G$.  Its trace is therefore the invariant dimension.  Characters
multiply on tensor products, while
\[
 \sum_{a\ge0}\Tr(\Sym^aA)t^a=\det(I-tA)^{-1}.
\]
Multiplying one copy for each variable and averaging proves
\eqref{eq:molien}.  The formal identity
\[
 -\log\det(I-tA)=\sum_{r\ge1}\frac{\Tr(A^r)}r t^r
\]
and $\Tr(\rho(g)^r)=\chi_V(g^r)$ prove \eqref{eq:power}; grouping elements by
class proves \eqref{eq:classpower}.  Finally
$h_a=\sum_{\lambda\vdash a}p_\lambda/z_\lambda$ gives
\eqref{eq:coefficient}.
\end{proof}

\begin{remark}[Character specializations]
For a genuine Steinberg or a specified genuine Weil representation, insert
its character into \eqref{eq:classpower}.  For a Deligne--Lusztig virtual
character $R_T^G(\theta)$ the same expression is interpreted in the completed
Grothendieck lambda-ring; its coefficients are virtual multiplicities and
need not be nonnegative.  These are specializations of the theorem, not
claims that all such character values collapse to one uniform product.
\end{remark}

\section{Permutation and parabolic quotient groups}

Suppose $G$ acts on a finite set $X$, and use the honest complex permutation
module $V=\C[X]$.  Let $c_r(g;X)$ count the $r$-cycles and put
$F_d(g;X)=|\Fix_X(g^d)|$.

\begin{proposition}[Permutation-cycle formula]\label{prop:permutation}
For every finite action,
\begin{align}
c_r(g;X)&=\frac1r\sum_{d\mid r}\mu(r/d)F_d(g;X),\label{eq:mobius}\\
\MGF_{G,X}&=\frac1{|G|}\sum_{g\in G}\prod_{i,r\ge1}
 (1-t_i^r)^{-c_r(g;X)},\label{eq:permutation}\\
[m_\tau]\MGF_{G,X}&=\#\left(G\backslash
 \prod_j\Sym^{\tau_j}X\right).\label{eq:orbits}
\end{align}
\end{proposition}

\begin{proof}
An $r$-cycle has determinant factor $1-t^r$ and contributes all of its $r$
points to $F_d$ precisely when $r\mid d$.  Hence
$F_d=\sum_{r\mid d}r c_r$, and M"obius inversion gives
\eqref{eq:mobius}.  The determinant factors give \eqref{eq:permutation}.
The monomial basis of $W_\tau(\C[X])$ is the displayed product of multiset
spaces, so its invariant dimension is the number of orbits, proving
\eqref{eq:orbits}.
\end{proof}

\begin{proposition}[Coset and parabolic cycles]\label{prop:coset}
For $X=G/H$,
\begin{align}
F_d(g;G/H)&=\frac{|C_G(g^d)|}{|H|}\,|(g^d)^G\cap H|,\label{eq:cosetfixed}\\
c_r(g;G/H)&=\frac1r\sum_{d\mid r}\mu(r/d)
 \frac{|C_G(g^d)|}{|H|}\,|(g^d)^G\cap H|.\label{eq:cosetcycles}
\end{align}
Substitution in \eqref{eq:permutation} gives the complete exact sigma-MGF.
In particular this applies to $G=\mathbf G^F$ and every $F$-stable parabolic
$P$, including projective points, Grassmannians, polar Grassmannians, and
fixed-shape partial flags.
\end{proposition}

\begin{proof}
The coset $xH$ is fixed by $g^d$ exactly when $x^{-1}g^dx\in H$.  Each element
of $(g^d)^G\cap H$ has $|C_G(g^d)|$ conjugators $x$.  Divide by $|H|$ for the
number of representatives of a coset and then apply \eqref{eq:mobius}.
\end{proof}

\section{Conjugation groups}

For conjugation on $G$ itself,
\[
 F_d(g;G_{\mathrm{conj}})=|C_G(g^d)|.
\]
Therefore
\begin{equation}\label{eq:conjcycles}
c_r^{\mathrm{conj}}(g)=\frac1r\sum_{d\mid r}\mu(r/d)|C_G(g^d)|,
\qquad
\MGF_{G,\mathrm{conj}}=\frac1{|G|}\sum_g\prod_{i,r}
(1-t_i^r)^{-c_r^{\mathrm{conj}}(g)}.
\end{equation}
On a fixed conjugacy class $C$, replace the fixed-point count by
$|C\cap C_G(g^d)|$.  Hence the conjugation series is determined by
centralizers, class intersections, and power maps.

\section{Fixed-field stability for bounded flags}

Fix $q$ and a flag shape
$\mathbf d=(d_1<\cdots<d_s)$, with $d=d_s$ independent of ambient rank.  Let
$X_n=\operatorname{Flag}_{\mathbf d}(\F_q^n)$.

\begin{theorem}[Type-$A$ bounded-support stability]\label{thm:typeA}
For every partition $\tau$ of total degree $D$,
\begin{equation}\label{eq:typeAstable}
[m_\tau]\MGF_{GL_n(q),X_n}\quad\text{is independent of $n$ for }n\ge dD.
\end{equation}
The same holds for $PGL_n(q)$.  For $SL_n(q)$ and $PSL_n(q)$ the safe range is
$n>dD$.  The stable coefficient is
\begin{equation}\label{eq:stablecoefficient}
a_{\tau,\mathbf d}(q)=\sum_{r=0}^{dD}\#\left(
GL_r(q)\backslash\mathcal C_{\tau,\mathbf d}^{\mathrm{span}}(r)\right),
\end{equation}
where the configuration space consists of $\tau$-colored flag multisets whose
top subspaces span $\F_q^r$.
\end{theorem}

\begin{proof}
A degree-$D$ configuration uses $D$ flags with multiplicity.  The sum of their
top subspaces, called the support, has dimension at most $dD$.  Two
configurations are $GL_n(q)$-conjugate exactly when their supported labelled
configurations are linearly isomorphic, because every isomorphism between the
supports extends across complements.  Once $n\ge dD$, every possible support
already embeds, so no new orbit type can appear.  Sorting by the support
dimension gives \eqref{eq:stablecoefficient}.  Scalar matrices act trivially,
so the $PGL$ orbits agree.  If $n>dD$, a nonzero complementary direction can
correct the determinant of any extension while fixing the target support;
this gives the stated $SL$ and $PSL$ range.
\end{proof}

\begin{proposition}[Formed-space towers]\label{prop:formed}
Fix one symplectic, unitary, odd-orthogonal, or one of the two
even-orthogonal towers, and let $X_n$ be the totally isotropic or singular
flags of fixed shape.  For every $\tau$ there is a finite threshold
$N_{\mathcal T}(\mathbf d,\tau)$ after which
$[m_\tau]\MGF_{G_n,X_n}$ is constant.
\end{proposition}

\begin{proof}
The span of a degree-$D$ configuration again has bounded dimension.  Its orbit
is determined by the labelled flags together with the restricted Hermitian,
alternating, or quadratic form and the multiplicities.  In characteristic two
the quadratic form, not merely its polar form, is retained.  There are only
finitely many such bounded-dimensional supported types over a fixed field.
Witt extension embeds and identifies each type once the chosen tower is large
enough; taking the maximum of these finitely many embedding ranks gives the
threshold.  Determinant or spinor-norm correction may require enlarging the
unused nondegenerate complement before passing to special, derived, or
projective groups.
\end{proof}

\begin{remark}[Why complete flags are excluded]
For $G/B$, Bruhat decomposition gives
$[m_{(1,1)}]\MGF_{G,G/B}=|B\backslash G/B|=|W|$.  This grows with rank, so
bounded shape is an essential hypothesis.
\end{remark}

\section{Representation scope for finite Lie-type groups}

Projective, subspace, flag, and conjugacy actions are honest complex
permutation representations.  Formula \eqref{eq:classpower} also applies to
every specified genuine complex character, including Steinberg and chosen
Weil representations.  The defining-characteristic action on
$\mathfrak g(\F_q)$ requires care:
\begin{itemize}
\item as a finite set with conjugation action it is an honest permutation
representation, with
$F_d(g)=q^{\dim_{\F_q}C_{\mathfrak g}(g^d)}$;
\item as a linear module over $\F_q$ it is not automatically a complex
representation, so a cross-characteristic realization or lift must be
specified before applying \eqref{eq:power}.
\end{itemize}

\section{Group-separated exact verification}

The executable companion is
\path{marimo_notebooks/power_character_stability.py}.
It calls a plain Python verification module, so the same checks run without a
notebook server.

\subsection{Linear group: $S_3\cong GL_2(2)$ and its Steinberg module}

Every one of the six integral $2\times2$ matrices of the standard $S_3$
representation is constructed.  For each $\tau$, the program builds the
matrices of every $\Sym^{\tau_j}V$, takes their Kronecker product, averages the
actual matrices, and computes the Reynolds rank.  Independently it uses only
$\chi(g^r)$ and Newton's recurrence to evaluate \eqref{eq:coefficient}.  In
the order
\[
(1),(2),(1,1),(3),(2,1),(1,1,1),
\]
both routes give
\[
(0,1,1,1,1,1).
\]
The largest projector-idempotence error is $5.56\times10^{-17}$.

\subsection{Parabolic group: $GL_3(2)$ on $G/P$}

The program enumerates all $168$ invertible $3\times3$ matrices over $\F_2$
and their action on the seven projective points.  It separately computes all
six conjugacy classes, the point stabilizer $P$ of order $24$, centralizers,
and class intersections with $P$.  Formula \eqref{eq:cosetfixed} reproduces
every needed fixed-point count and reconstructs the cycles of all $168$
elements.  Actual permutation-matrix averages, cycle factors, and direct
multiset-orbit traversal agree on
\[
(1,2,2,4,5,6).
\]
There are $715$ direct fixed-point equalities and $168$ full cycle
reconstructions.

\subsection{Conjugation group: $S_3$ on itself}

The six conjugation permutations are built from the group law.  Direct fixed
points agree with $|C_G(g^d)|$ for every required power, and all six cycle
decompositions are reconstructed.  Matrix averages, \eqref{eq:conjcycles},
and direct orbits agree on
\[
(3,8,11,17,32,49).
\]

\subsection{Stable type-$A$ groups}

Transvections generate the actions without enumerating the large ambient
groups.  Direct configuration-orbit traversal gives:
\begin{center}
\begin{tabular}{@{}lrrrr@{}}
\toprule
$\tau$ for projective points & $n=1$ & $n=2$ & $n=3$ & $n=4$\\
\midrule
$(1)$       &1&1&1&1\\
$(2)$       &1&2&2&2\\
$(1,1)$     &1&2&2&2\\
$(3)$       &1&3&4&4\\
$(2,1)$     &1&4&5&5\\
$(1,1,1)$   &1&5&6&6\\
\bottomrule
\end{tabular}
\end{center}
The onset is exactly consistent with $n\ge|\tau|$ when $d=1$.  For
$GL_n(2)$ on $\operatorname{Gr}_2(n)$, both $\tau=(2)$ and $(1,1)$ give three
orbits at $n=4$ and again at $n=5$, matching the threshold $n\ge2|\tau|=4$
and the intersection-dimension classification.

\resultbox{\textbf{Verification result.} All 26 reported coefficient or
stability rows pass.  These computations test the formulas at small
representative dimensions with distinct failure modes; the all-rank claims
rest on the proofs above.}

\section{Reproduction}

From the repository root, run
\begin{verbatim}
python3 -m for_this_guy.power_character_bounded_flag_stability.verification
.venv/bin/marimo run marimo_notebooks/power_character_stability.py
\end{verbatim}

\begin{thebibliography}{9}
\bibitem{Howe}
S. Howe, \emph{Random matrix statistics and zeroes of $L$-functions via
probability in lambda-rings}, \href{https://arxiv.org/abs/2412.19295}{arXiv:2412.19295}.

\bibitem{Huang}
J. Huang, \emph{Some extensions of Witt's theorem},
\href{https://arxiv.org/abs/0705.3839}{arXiv:0705.3839}.

\bibitem{Thomas}
T. Thomas, \emph{The character of the Weil representation},
\href{https://arxiv.org/abs/math/0610644}{arXiv:math/0610644}.

\bibitem{MWZ}
P. Magyar, J. Weyman, and A. Zelevinsky, \emph{Multiple flag varieties of
finite type}, \href{https://arxiv.org/abs/math/9805067}{arXiv:math/9805067}.
\end{thebibliography}

\end{document}
